\documentclass[11pt]{article}
\usepackage[margin=1in]{geometry}
\usepackage{amsmath,amssymb,amsthm}
\usepackage{graphicx}
\usepackage{hyperref}
\usepackage{xcolor}

\newtheorem{theorem}{Theorem}
\newtheorem{lemma}[theorem]{Lemma}
\newcommand{\C}{\mathbb C}

\title{A Literal Counterexample to the Displayed\\
Daws Cotype Conjecture in arXiv:1901.07496}
\author{Automatic mathematics run \texttt{fa\_banach\_001}}
\date{13 August 2026}

\begin{document}
\setlength{\emergencystretch}{3em}
\maketitle

\begin{center}
\fbox{\parbox{0.91\textwidth}{
\textbf{Status.} Candidate counterexample, likely valid, to Conjecture~2
\emph{as literally displayed} in Gerasimova--Thom.  The cited Daws
conjecture uses a corrected exponent, so this packet does not claim to
disprove Daws's intended conjecture or Gerasimova--Thom Conjecture~1.
}}
\end{center}

\section{Source conjecture}

The source is M. Gerasimova and A. Thom, \emph{On the isometrisability of
group actions on $p$-spaces}, arXiv:1901.07496v2 \cite{source}.  On printed
page~7 it defines $B_p(\Gamma)$ using coefficients of isometric
representations on $q$-spaces and then displays:
\begin{quote}
\textbf{Conjecture 2 (Daws).} The space $B_p(\Gamma)$ has cotype $q$,
where $\frac1p+\frac1q=1$.
\end{quote}
No range for $p$ is included in the displayed statement.  The complete
source context is shown below.

\begin{center}
\includegraphics[width=0.98\textwidth]{figures/open_problem_crop.png}
\end{center}

\section{Definitions}

A Banach space $X$ has (Rademacher) cotype $r$ with constant $C$ if, for
every finite family $x_1,\ldots,x_n\in X$,
\[
 \left(\sum_{j=1}^n\|x_j\|^r\right)^{1/r}
 \le C\left(\mathbb E_\varepsilon
       \left\|\sum_{j=1}^n\varepsilon_jx_j\right\|^2\right)^{1/2},
\]
where the $\varepsilon_j$ are independent random signs.  Standard
terminology only assigns cotype exponents in $[2,\infty)$; the displayed
inequality also makes the obstruction below meaningful for $r<2$.

The source's coefficient space $B_p(\Gamma)$ contains the coefficient of
every isometric representation on a $q$-space, with $q=p/(p-1)$.  In
particular the trivial one-dimensional representation on $\C$ is allowed.

\section{Idea of the proof}

The obstruction is independent of group representation theory.  Repeating
one nonzero vector $n$ times makes the left side of the cotype inequality
grow as $n^{1/r}$, whereas the Rademacher square mean on the right grows only
as $n^{1/2}$.  Thus no nonzero Banach space can have cotype below~$2$.
For $p>2$ the conjugate exponent in the displayed conjecture is below~$2$,
and $B_p(\Gamma)$ is nonzero because it contains the constant coefficient~1.

\section{Counterexample}

\begin{lemma}\label{lem:no-small-cotype}
No nonzero Banach space has Rademacher cotype $r<2$.
\end{lemma}

\begin{proof}
Let $0\ne x\in X$ and apply the cotype inequality to
$x_1=\cdots=x_n=x$.  Independence and
$\mathbb E|\sum_{j=1}^n\varepsilon_j|^2=n$ give
\[
 n^{1/r}\|x\|
 \le C\sqrt n\,\|x\|.
\]
Hence $n^{1/r-1/2}\le C$ for every $n$, impossible when $r<2$.
\end{proof}

\begin{theorem}\label{thm:counterexample}
Conjecture~2 of arXiv:1901.07496v2 is false as literally displayed.
Explicitly, take
\[
 p=3,\qquad q=\frac32,\qquad \Gamma=\{e\}.
\]
Then $B_3(\Gamma)\cong\C$ is nonzero and does not have cotype $3/2$.
More generally, the displayed conclusion fails for every $2<p<\infty$ and
every discrete group $\Gamma$.
\end{theorem}

\begin{proof}
The trivial representation of $\Gamma$ on the one-dimensional $q$-space
$\C$ is isometric, and its coefficient obtained from $x=1$ and the identity
functional is the constant function~$1$.  Thus $B_p(\Gamma)\ne\{0\}$.
If $p>2$, its conjugate exponent $q=p/(p-1)$ satisfies $q<2$, so
Lemma~\ref{lem:no-small-cotype} rules out cotype~$q$.  For the trivial group,
every coefficient is a scalar function on one point, hence
$B_3(\{e\})\cong\C$.
\end{proof}

\section{Convention diagnosis and limitations}

The cited paper of Daws does \emph{not} state the unrestricted formula in
the displayed source conjecture.  In Section~8.3, immediately after
Proposition~8.8, Daws writes that Runde's algebra $B_p(G)$ is conjectured to
have cotype
\[
 \max(p,p')
\]
\cite[p.~74]{daws}.  Moreover, Daws explicitly warns that his Herz convention
sometimes interchanges $p$ and $p'$ relative to Runde's convention.  These
two details explain the mismatch.  The exponent $\max(p,p')$ is always at
least~$2$ and survives the elementary obstruction above.

Consequently, Theorem~\ref{thm:counterexample} resolves only the literal
unrestricted displayed Conjecture~2.  It does not settle the corrected Daws
conjecture, and it gives no counterexample to Conjecture~1 concerning
$(2,p)$-isometrisability and Littlewood functions.

\section{Verification and novelty check}

The proof uses only the defining cotype inequality and the trivial
one-dimensional coefficient; there is no numerical or computer-assisted
step.  A separate proof audit is included as \texttt{verification.md}.

A bounded search through 13 August 2026 covered the run's four lightweight
indexes, the current source and bibliography, the exact displayed phrase,
the title and citations of arXiv:1901.07496, and arXiv-oriented searches for
$B_p(G)$ with cotype.  No later paper explicitly flagging this displayed
range error was found.  The original cited Daws paper was inspected directly
and supplies the corrected $\max(p,p')$ formulation.  Novelty confidence is
therefore limited to the explicit diagnosis and counterexample to the source
statement; the intended corrected conjecture is known from Daws and remains
open.

\section{Human-review recommendation}

The mathematical counterexample is elementary.  Review should focus on
editorial scope: whether the displayed statement is meant to inherit an
unstated range or Daws's alternate convention.  The packet should remain
labelled ``literal counterexample'' and must not be reported as a disproof of
the corrected Daws conjecture.

\begin{thebibliography}{9}
\bibitem{source}
M. Gerasimova and A. Thom,
\emph{On the isometrisability of group actions on $p$-spaces},
arXiv:1901.07496v2, 2020.

\bibitem{daws}
M. Daws,
\emph{$p$-operator spaces and Fig\`a--Talamanca--Herz algebras},
Journal of Operator Theory 63 (2010), 47--83;
arXiv:math/0611833.
\end{thebibliography}

\end{document}

